\begin{abstract}
Сүүлийн жилүүдэд хүний зан үйл, дадал хэвшлийг өөрчлөхөд дижитал технологийг ашиглах асуудал нь зан үйлийн шинжлэх ухаан болон програм хангамжийн инженерийн огтлолцолд байрлах чухал судалгааны чиглэл болж өргөжиж байна. Эрүүл мэнд, боловсрол, бүтээмж, хувийн зохион байгуулалт зэрэг олон талбарт давтагдах зан үйлийг тогтвортой болгох хэрэгцээ өндөр боловч түгээмэл дадал хянах системүүдийн дийлэнх нь сануулга, тасралтгүй гүйцэтгэлийн цуваа, гүйцэтгэлийн бүртгэл зэрэг хэсэгчилсэн механизмд төвлөрч, дадал үүсэх онолын суурь механизмыг системийн түвшинд бүрэн тусгаж чаддаггүй.

Энэхүү судалгааны ажил нь дадал төлөвшлийн мөчлөг, Фоггийн зан үйлийн загвар, итгүүлэн нөлөөлөх системийн зохиомжийн зарчим, SRBAI-д суурилсан автоматжилтын хэмжилтийг нэгтгэн зан үйлд суурилсан дадал төлөвшүүлэх програм хангамжийн системийн загвар боловсруулах зорилготой. Судалгааны хүрээнд өдөөгч дохио, хүсэл тэмүүлэл, хариу үйлдэл, урамшуулал гэсэн бүрэлдэхүүнүүдийг системийн шаардлага, архитектур, өгөгдлийн загвар, үндсэн алгоритм, дасан зохицох сануулгын логиктой уялдуулан загварчилсан. Мөн дадлын хүчийг зөвхөн гүйцэтгэлийн тоо эсвэл цуваагаар бус, автоматжилтын оноо, гүйцэтгэлийн тогтвортой байдал, нөхцөл байдлын тогтвортой байдал зэрэг үзүүлэлтүүдийг нэгтгэсэн үйл ажиллагааны загвараар үнэлэх аргыг санал болгосон.

Судалгааны үр дүнд зан үйлд суурилсан системийн үзэл баримтлалын загвар, функциональ болон функциональ бус шаардлага, хэрэглэгч--сервер хэв маягт суурилсан модульчилсан архитектур, өгөгдлийн хамаарлын загвар, урьдчилсан үнэлгээний арга зүйг тодорхойлов. Энэхүү ажил нь зан үйлийн онолыг програм хангамжийн хэрэгжихүйц шийдэл болгон хөрвүүлэх онол-практикийн суурь загвар боловсруулахад чиглэж байгаа бөгөөд цаашдын туршилтын загварын хэрэгжүүлэлт, өргөтгөсөн туршилт, урт хугацааны үнэлгээний ажлын үндэс болно.
\end{abstract}
